\section{Discussion, limitations, and extensions}

The lower-bound analysis isolates one obstruction that is already present in the standard interior partial mean problem. The target in \cref{obj:def:theta-functional} is an observed-data partial mean at a fixed interior dose, and the risk in \cref{obj:def:minimax-risk} is taken over the same Hölder dose class of \cref{obj:def:holder-dose-class}. Within that class, \cref{obj:thm:sharp-pointwise-lower-bound} gives a treatment-regression lower floor. The comparison statements in \cref{obj:thm:sharp-minimax-smooth-covariate} and~\cref{obj:thm:frontier-bracket-deficient} relate this floor to the published benchmark sequence \(\rho_n\) from \cref{obj:def:published-hoif-rate}, while \cref{obj:lem:published-upper-bound-cited} supplies the external Bonvini--Kennedy comparator on its localized-regularity class.

For each fixed \(\beta>0\), under the corresponding slack-baseline condition in \cref{obj:ass:baseline-submodel-slack}, the certified exponent of the same-class lower bound is invariant to the treatment-density smoothness index. The constant in the lower bound, and the feasibility of the slack baseline itself, may depend on \(\beta\) and on the other fixed model constants. Thus the conclusion is exponent-level insensitivity to \(\beta\), with constants calibrated separately for each fixed smoothness value.

This distinction matters most when \(4s<d\). In that regime, \cref{obj:thm:frontier-bracket-deficient} shows that the published benchmark sequence has a smaller exponent than the same-class treatment-regression lower floor. The unresolved same-class question is whether the upper behavior in low covariate smoothness follows the published benchmark, an intermediate exponent, or a rate that depends more directly on the smoothness of the treatment density.

The scope is tied to the interior positivity regime. \Cref{obj:ass:local-positivity} imposes a fixed lower bound on the treatment density in a neighborhood of \(t_0\), and \cref{obj:ass:interior-dose} selects an interior evaluation window. These restrictions match the usual interior continuous-exposure formulation. Weak-overlap, design-adaptive, and trimmed targets lead to a different statistical experiment in which the amount of information near \(t_0\) enters the rate calculation directly. Recent work on continuous exposures and overlap-sensitive causal estimands emphasizes these complementary regimes \citep{hejazi2022,doss2022,hudson2023,shi2024,zhang2024,schindl2024,bonvini2024}.

Neighboring estimands raise a similar qualification. The partial mean studied here evaluates the regression surface at a single interior dose and averages over the marginal law of \(X\). Other continuous-exposure targets average over dose neighborhoods, consider stochastic interventions, or target modified exposure policies. Their definitions can smooth the treatment coordinate or alter the role of the conditional treatment density. The present lower bound is consequently a benchmark for the fixed-dose interior partial mean, while those neighboring functionals invite estimand-specific minimax analyses.

Finally, the comparison with higher-order influence-function methods is a comparison across stated model classes. The imported result in \cref{obj:lem:published-upper-bound-cited} supplies an upper-side benchmark under additional localized regularity, while the lower bound here is proved over the Hölder class in \cref{obj:def:holder-dose-class}. Two routes can bridge the statements: a same-class upper bound under \cref{obj:def:holder-dose-class}, or an inclusion argument covering the present Hölder class by the localized-regularity class with the needed constants.
